\documentclass[11pt,a4paper]{article}
\usepackage{fontspec,amsmath,amssymb,amsthm,microtype}
\setmainfont{lmroman10-regular.otf}[BoldFont=lmroman10-bold.otf,ItalicFont=lmroman10-italic.otf,BoldItalicFont=lmroman10-bolditalic.otf]
\usepackage[margin=25mm]{geometry}
\usepackage[unicode,hidelinks]{hyperref}
\hypersetup{pdfauthor={LangevinLean Project}}
\setlength{\emergencystretch}{3em}
\setlength{\parskip}{0.4em}
\setlength{\parindent}{0pt}
\allowdisplaybreaks[2]
\raggedbottom
\title{\texorpdfstring{Differentiation of the entropy integral}{Differentiation of the entropy integral}}
\author{LangevinLean Project}
\date{}
\begin{document}\maketitle
Draft mathematical exposition; these notes may contain mistakes.


For a positive differentiable function \(r\), the scalar derivative is \((r\log r)'=r'(\log r+1)\). Applying this identity to a density depending on time gives a derivative of its entropy integral when the differentiated integrands have a common integrable bound on a time neighborhood. This bound and the relevant measurability and integrability assumptions are hypotheses of the result.


\section{\texorpdfstring{The derivative of a positive function times its logarithm}{The derivative of a positive function times its logarithm}}


Assume \(r\) is differentiable at \(t\) and \(r(t)>0\). The product rule and the derivative of the logarithm give \((r\log r)'(t)=r'(t)\log r(t)+r(t)r'(t)/r(t)\). Positivity implies \(r(t)\ne0\), so cancellation yields \(r'(t)(\log r(t)+1)\). The constant term is retained at this stage. Its contribution can be removed from an entropy calculation only after the derivative of total mass has been shown to vanish.


\section{\texorpdfstring{Differentiation under the integral sign}{Differentiation under the integral sign}}


Fix \(t\), and assume that almost every spatial function \(u\mapsto r_u(x)\) is positive and differentiable on a common neighborhood of \(t\). Suppose there is an integrable function \(M(x)\) such that \(|\partial_u r_u(x)(\log r_u(x)+1)|\le M(x)\) on that neighborhood. Assume also measurability of the integrands and integrability of \(r_t\log r_t\). The scalar chain rule identifies the pointwise derivative, and the dominated differentiation theorem gives

\[\frac{d}{dt}\int r_t\log r_t\,d\nu=\int\dot r_t(\log r_t+1)\,d\nu.\]

The same bound proves integrability of the derivative integrand. Uniformity on a time neighborhood, rather than a bound at one time alone, is what justifies exchanging differentiation and integration.


\end{document}
